\section{Introduction}
Modern language models acquire knowledge and capabilities through
pre-training and post-training. After deployment, however, their
parameters are typically frozen. The behavior observed at inference
time therefore depends not only on what has been learned by the model,
but also on the context dynamically constructed around each inference.

A simplified view is

\[
    y_t = f_{\theta}(C_t),
\]

where $\theta$ represents the frozen model parameters and $C_t$
represents the context available at inference time.

This project studies the construction of $C_t$ as a systems problem.
Rather than training a new model, the objective is to investigate how
a frozen model can be provided with an effective working context under
a limited context budget.

\subsection{From Model capabilities to effective behavior}
Information learned by a language model should not be viewed simply as
an unstructured information store. Training produces distributed
representations and computational capabilities that depend on the
training data, model architecture, training objectives, and
post-training process.

However, the existence of a capability in the model does not guarantee
that it will be expressed in every inference.

Conceptually,

\[
    \text{Learned Capability}
    +
    \text{Appropriate Context}
    \rightarrow
    \text{Effective Behavior}.
\]

This distinction between \emph{having} a capability and
\emph{eliciting} that capability motivates the study of context
engineering.


\subsection{Context as an Inference-Time Program}

A useful perspective is to view context as an inference-time program
for a frozen model.

The analogy is not exact, but the frozen model can be regarded as a
learned computational system, while the context provides the
instructions, data, state, constraints, and available resources for a
particular execution.

A modern context may contain

\[
C_t =
\{
    \text{instructions},
    \text{history},
    \text{memory},
    \text{evidence},
    \text{tools},
    \text{observations},
    \text{state}
\}.
\]
From this perspective, prompt engineering is only one part of a broader
context-engineering problem.

Context engineering serves two related purposes:

\[
\boxed{
    \text{Capability Elicitation}
    +
    \text{Runtime Harness}
}
\]

It provides information and capabilities to the model while also
guiding and constraining how the model interacts with its environment.

\subsection{Context Engineering in Agentic Systems}
The problem becomes more important in agentic systems because an agent
continuously accumulates information while interacting with its
environment.

A simplified agent loop can be represented as

\[
    \text{Observe}
    \rightarrow
    \text{Reason}
    \rightarrow
    \text{Act}
    \rightarrow
    \text{Observe}.
\]

During long-running execution, conversation history, retrieved
documents, source code, tool outputs, execution traces, errors, and
previous decisions continuously accumulate.

Let $H_t$ represent all information available to the system. In a
long-running agent,

\[
    |H_t| \gg |C_t|.
\]

The agent harness must therefore construct a smaller working context:

\[
    H_t
    \xrightarrow{\pi}
    C_t
    \xrightarrow{f_{\theta}}
    y_t,
\]

where $\pi$ represents a context-construction policy.

This problem is already visible in modern coding agents and agentic
applications, where repositories, terminal outputs, tool results,
documentation, and interaction histories compete for limited context.

Furthermore, different workloads may require different context
structures:

\[
    C_{\text{chat}}
    \neq
    C_{\text{reasoning}}
    \neq
    C_{\text{coding}}
    \neq
    C_{\text{multimodal}}.
\]

This suggests that a single static context-management strategy may not
be optimal for every model, task, and system state.


\section{Research Gap}

Existing research demonstrates that long contexts may be difficult to
use effectively and that explicit memory management can improve
long-running AI systems.

However, several questions remain relevant as models and agent
architectures continue to evolve.

First, it is unclear whether the U-shaped
\emph{Lost-in-the-Middle} \cite{lost-in-the-middle} \cite{hsieh-etal-2024-found} \cite{zhang2024found} effect remains significant in current
generation language models and how this effect changes with context
length.

Second, common context-management strategies such as full history,
fixed windows, summarization, and relevance-based retrieval are often
studied independently. Their quality--cost trade-offs under controlled
conditions require further comparison.

Third, different models, tasks, and agent states may require different
working contexts. This motivates investigating whether the context
strategy itself should be adaptive.

The proposed research therefore asks:

\begin{enumerate}

    \item Does the Lost-in-the-Middle effect still exist in modern
    language models, and how does it change as context length increases?

    \item Can selective context strategies maintain or improve task
    quality compared with providing the complete available context?

    \item Can a lightweight adaptive policy select an appropriate
    context strategy according to the current task and system state?

\end{enumerate}

The first experiment will vary both relevant-information position and
total context length:

\[
    Q = f(p, L),
\]

where $p$ represents the relative position of relevant information and
$L$ represents total context length.

Subsequent experiments will compare several context strategies:

\[
\begin{aligned}
    &\text{Full History},\\
    &\text{Fixed Window},\\
    &\text{Summarization},\\
    &\text{Relevance Retrieval},\\
    &\text{Adaptive / Hybrid Selection}.
\end{aligned}
\]

The project will initially keep the model parameters frozen and will
not use reinforcement learning to train a separate manager model.
Instead, the emphasis will be on implementation, controlled
experiments, and interpretable context-selection policies.

The overall research problem can be summarized as

\[
\boxed{
    C_t^{*}
    =
    \pi(
        \text{model},
        \text{task},
        \text{state},
        \text{history},
        \text{budget}
    )
}
\]

with the objective of finding a better trade-off between task quality
and context cost.

The central hypothesis is therefore not simply that shorter context is
better, but that

\[
\boxed{
    \text{More Context}
    \not\equiv
    \text{Better Context}.
}
\]

The broader question is:

\begin{quote}
How should an agentic system construct the right working context for
the right model, task, and state?
\end{quote}

\section{Research Questions}
The proposed research investigates the following questions:

\begin{enumerate}
    \item \textbf{RQ1}
    \begin{quote}
        How does the position of relevant information within the context affect its utilization by modern LLMs?
    \end{quote}
    \item \textbf{RQ2}
    \begin{quote}
        How does the length of the context influence the relationship between information position and model performance?
    \end{quote}
    \item \textbf{RQ3}
    \begin{quote}
        How does information interference from semantically similar distractors affect the utilization of relevant information at different positions?
    \end{quote}
\end{enumerate}
